\subsection{Exact block cuts and sixteen-arrow star reductions}

\begin{corollary}[Height-one support]
\label{C:height-one-zero-H2-products}
Let $C$ be a finite category with no composable pair of arrows
between distinct objects.  If every corner has zero rational
second cohomology, all products in $H^2(BC;\mathbb Q)$ vanish.
\end{corollary}
\begin{proof}
Partition the objects into sources and sinks, assigning isolated
objects arbitrarily.  Each full block is a disjoint union of
corner monoids, so its second cohomology is zero.  No arrow
goes from the sink block to the source block.
Lemma~\ref{C:two-block-zero-H2-products} applies when both
blocks are nonempty.  If one is empty, the category itself is
a disjoint union of corners and has zero second cohomology.
\end{proof}

\begin{proposition}[Sufficient zero-cohomology flags for the search]
\label{C:exact-block-cut-flags}
The following are sufficient conditions for a full subcategory
$D$ to have $H^2(BD;\mathbb Q)=0$: each connected component
satisfies at least one of the following conditions.
\begin{enumerate}
\item it is a single corner whose normalized second homology
over $\mathbb F_2$ is zero, or whose rational algebra is semisimple;
\item it has two objects and at most five arrows, or exactly
one cross arrow, or is triangular with both corner algebras
semisimple over $\mathbb Q$;
\item it is connected, has three objects, and at most seven arrows.
\end{enumerate}
Thus an oriented partition of the objects into two blocks satisfying
these conditions excludes a projective-plane square, provided no
arrow goes from the second block to the first.
\end{proposition}
\begin{proof}
Cohomology of a finite disjoint union is the direct sum in degree
two.  For a corner of order $q$, the normalized bar groups in
degrees one through three have dimensions $q-1$, $(q-1)^2$,
and $(q-1)^3$.  Over $\mathbb F_2$ the two relevant differentials are
\[
 \partial[a|b]=[b]+[ab]+[a],\qquad
 \partial[a|b|c]=[b|c]+[ab|c]+[a|bc]+[a|b],
\]
with any identity-containing term omitted.
Hence the exact dimension is
$(q-1)^2-\operatorname{rank}\partial_2-
\operatorname{rank}\partial_3$.  An integral matrix has rational
rank at least its rank modulo two.  A zero value therefore forces
zero second rational homology, and then zero rational cohomology.
For a semisimple rational corner algebra all positive Ext groups
vanish.

Lemma~\ref{C:two-object-five-H2-zero} proves the five-arrow case.
A singleton cross hom-set gives the join of two connected corner
spaces, which has zero second rational cohomology.  In the
triangular semisimple case the relative bar over the product of
the two corner algebras has no term of length two, and exact
corner coinvariants leave no second cohomology.
Lemma~\ref{C:three-small-H2-zero} proves the three-object case.
The last assertion is precisely
Lemma~\ref{C:two-block-zero-H2-products}.

The implementation \path{code/c_cp2_block_prunes.hpp} computes
the two displayed ranks by exact binary elimination.  Its uses
have $q\leq8$, so the largest binary row has $(q-1)^2\leq49$
positions and fits in the stated 64-bit word.  It first splits
each proposed block into its actual connected components, and
tests every nonempty proper object subset for the absence of
reverse arrows.  A failed sufficient test never removes a case.
\end{proof}

\begin{proposition}[Four-object stars through sixteen arrows]
\label{C:sixteen-four-star-prunes}
The three exclusions in
Proposition~\ref{C:fifteen-four-star-prunes} remain valid for a
four-object inward or outward star with at most sixteen arrows.
\end{proposition}
\begin{proof}
The first exclusion uses only zero second cohomology at the
three leaves and a singleton incoming arrow fixed by the central
monoid.  Its proof by Lemma~\ref{C:fixed-incoming-cups} has no
fifteen-arrow restriction.

For a central monoid of order at most nine and leaves of orders
at most two, the exact leaf-coinvariant functors in the earlier
proof still give a rational cohomology-ring equivalence.  Identifying
the resulting three sink points gives the ordinary act cone with
two additional circle summands.  A hypothetical projective plane
would give bounded rational homology, degree-two rank one, and
a nonzero square, contrary to
Lemma~\ref{C:nine-bounded-rank-one}.  Its hypotheses explicitly
allow the additional first homology.

Finally two singleton legs give the same actual splitting
\eqref{C:four-star-two-singleton-splitting}.  The primitive
projective-plane square must lie entirely in one wedge summand.
The suspension summand has zero products, and consequently must
be integrally acyclic.  Its desuspension is the central monoid
space; Theorem~\ref{A:bounded} and the homological Whitehead
theorem make that space contractible.  The remaining two joins
are contractible, and the whole category reduces to its
center--third-leaf collage.  That collage has at most
\[
 16-(|A_1|+1)-(|A_2|+1)\leq12
\]
arrows, already excluded by Theorem~\ref{C:CP2-lower16}.
All three arguments are therefore valid at sixteen arrows.
\end{proof}
